\section{Conclusion}\label{sec:conclusion}

\textbf{RQ1.} OR research on last-mile delivery has grown by \BibGrowthPct\% a year since 2015, led by vehicle routing with drones, robots and two-echelon systems solved by mixed-integer programming and large neighbourhood search; quick commerce (\BibQComm{} of \BibIncluded{} articles), India (none) and rider welfare (\BibWelfare) remain almost unstudied, and the delivery promise is rarely a decision variable.
\textbf{RQ2.} In Bengaluru the cost-optimal network has about \StoresTFifteen{} dark stores for any promise of 15 minutes or more; a 10-minute promise cannot reach \UnreachTenPct\% of residents from mapped retail sites, and serving everyone within 10 minutes needs \GridStorePremiumPct\% more stores at \GridCostPremiumPct\% higher cost per order. Rider speed, not the promise, drives cost: a 15~km/h cap costs about \SpeedCapPremiumPct\% more per order than 20~km/h.
\textbf{RQ3.} A promise is realistic where an address lies within the safe ride radius of a store; on that basis 15 minutes is realistic for nearly all of Bengaluru and 10 minutes for \MixedTenShare\% of residents, and a promise differentiated this way costs only \MixedPremiumPct\% more per order than a uniform 15-minute network. Regulating speed norms and disclosure of delivery radii is more meaningful than regulating advertised minutes.
